% Generated draft: requires mathematical and source review.
\documentclass[11pt]{article}
\usepackage[margin=25mm]{geometry}
\usepackage{fontspec}
\setmainfont{Latin Modern Roman}
\usepackage{amsmath,amssymb,mathtools,unicode-math}
\setmathfont{Latin Modern Math}
\usepackage{enumitem,needspace,xurl,hyperref}
\setlength{\parindent}{0pt}
\setlength{\parskip}{5pt}
\setlength{\emergencystretch}{3em}
\newcommand{\N}{\mathbb{N}}
\newcommand{\Z}{\mathbb{Z}}
\newcommand{\Q}{\mathbb{Q}}
\newcommand{\R}{\mathbb{R}}
\newcommand{\C}{\mathbb{C}}
\newcommand{\F}{\mathbb{F}}
\newcommand{\A}{\mathbb{A}}
\newcommand{\PP}{\mathbb{P}}
\newcommand{\E}{\mathbb{E}}
\newcommand{\eps}{\varepsilon}
\newcommand{\dd}{\,\mathrm{d}}
\begin{document}
\section*{UnsolvedMath ID: 30002006}
\textbf{Source problem code:} catalogue entry
\subsection{Definitions and mathematical statement}
Let \(A_{\mathrm{crit}}\) denote the critical-level affine enveloping algebra intended by the source, let \(Z_{\mathrm{crit}}\) denote its relevant center, and let \(\mathcal M\) denote the intended abelian category of modules. These are provisional symbols: the supplied record does not specify the underlying affine Lie algebra, ground field, completion, or continuity conditions on modules and central actions. The statement refers simply to modules, whereas the supplied background identifies the intended target as the full smooth critical category; these descriptions cannot be identified without further hypotheses.

Under the usual categorical interpretation, a block is a nonzero abelian category that is indecomposable: it is not equivalent to a product of two nonzero abelian categories. A direct-product decomposition into blocks means an equivalence of abelian categories
\[
\Phi:\mathcal M\xrightarrow{\sim}\prod_{i\in I}\mathcal B_i,
\]
where \(I\) is a set and every \(\mathcal B_i\) is a block. Objects and morphisms in the product are families, with composition defined componentwise.

\textbf{Problem.} Using the action of \(Z_{\mathrm{crit}}\), determine whether the intended category \(\mathcal M\) admits a canonical direct-product decomposition into blocks as above.

\textbf{Unresolved qualifications.} The record does not specify how the center is to determine the indexing set or factors, what notion of canonicity is required, or whether the intended meaning of block agrees with the categorical interpretation given here. In particular, no decomposition by central characters, generalized central characters, or central idempotents is stipulated. The category and these requirements must be clarified before this becomes a uniquely specified mathematical assertion.
\subsection{Short English statement}
The problem asks whether the center at critical level canonically determines a product decomposition of the intended module category into indecomposable abelian categories. The supplied background specifies the full smooth category, but the statement itself leaves the module category and the meaning of canonicity unspecified. The supplied open status is dated curation, not an independently verified conclusion.
\subsection{Sources}
\begin{itemize}
\item Supplied local UnsolvedMath registry record, permanent ID \(30002006\), title \emph{Block Decompositions of Critical-Level Affine Modules}, source problem code \texttt{catalogue entry}, in the Oberwolfach Reports Open Problems collection. The record gives the proposed year as 2012 and names no individual proposer.
\item \emph{Enveloping Algebras and Geometric Representation Theory} (2012), Oberwolfach Reports, Report 13, pp.\ 733--809, DOI \texttt{10.4171/owr/2012/13}. Supplied report provenance and literature reference: \url{https://doi.org/10.4171/owr/2012/13}.
\item Supplied literature reference: \emph{Centers for the Restricted Category O}, \url{https://arxiv.org/abs/1301.6016}.
\item Supplied literature reference: \emph{D-modules on the Affine Flag Variety}, \url{https://arxiv.org/abs/0712.0788}.
\item Literature assessment dated exactly \texttt{2026-08-21}: the supplied curation labels the problem open and reports that centers and block or linkage decompositions are understood for restricted or localized critical categories, while the requested direct-product description of the full smooth critical category remains open. This is supplied dated curation, not independently verified literature assessment. The record expressly requests independent verification of time-sensitive status and cited claims.
\end{itemize}
\end{document}
